\section*{The Propers: Detailed Commentary}

This is the synthesis edition. In place of the full edition's complete
appointed texts and its element-by-element sweep, the commentary here is
organised by five claims that no single proper establishes and that only the
whole formulary, read against its checked reception, makes visible. Every
appointed element informs what follows; none has a heading of its own. The
complete appointed Latin and English, the collation against two 1962 editions
and against the Clementine Vulgate, and the individual dossiers of every
witness cited here are in the full edition and in the leaf's source records ---
\texttt{propers/verified.md} and \texttt{research/scope.md}.

\subsection*{1. Two texts of this Mass are the tradition's own language, promoted}

The ordinary relation between exegesis and liturgy is that the liturgy quotes
Scripture and the commentators explain it. Two elements of this formulary
invert or short-circuit that relation, in opposite directions, and both are
verifiable.

The Introit sings \latin{Deus qui inhabitáre facit unánimes in domo}. No Bible
of the Latin west reads \latin{unánimes} there: the Clementine has
\latin{unius moris}, and the Douay accordingly has ``of one manner,'' which is
why no registered English can be quoted for the sung word. What
\latin{unánimes} is, in the checked corpus, is a \textit{gloss}. Augustine's
lemma is \latin{unius modi} and he explains it on the spot ---
\latin{unanimes, unum sentientes} --- observing besides that the Greek will
bear either \latin{modi} or \latin{mores}. Peter Lombard transmits the
explanation with the tag attached. Bellarmine reads it back into the verse as
though it were the verse. Only Cassiodorus lemmatises \latin{unanimes} as text,
and he is the one witness who also restores the Acts citation Augustine had
dropped, identifying the house as the Church. The Roman chant sings the gloss.
This guide asserts no line of descent --- no \work{Psalterium Romanum} witness
was collated --- but the effect is that the assembly's first words on this
Sunday are a patristic reading that has hardened into a liturgical text.

The Gospel runs the traffic the other way. \latin{Ephphetha} left Mark and
became a rite. St.~Ambrose is the earliest witness and states the derivation
outright: ``celebrating the mystery of the opening, we said, Epphatha, which is,
Be opened \dots\ Christ made use of this mystery in the Gospel, as we read, when
He healed him who was deaf and dumb.'' \work{De sacramentis} adds the gesture
--- the priest touching ears and nostrils --- and explains the transfer from the
mouth on grounds of propriety. Fifteen centuries later the \work{Rituale
Romanum} of 1872 still prints \latin{Ephpheta, quod est, Adaperire} with the
same saliva and the same ears, at the baptism of infants and again at the
baptism of adults. Between the two lies an exegetical tradition that read the
spittle as divine wisdom (Bede) and as proof that every part of Christ's flesh
was holy (Theophylact) --- but the rite does not depend on that tradition; it is
older than most of it and travels beside it. A reader is likelier to have met
the word at a font than in Mark, and the missal's own rubric at the end of this
Gospel, \latin{Credo}, sits oddly close to the seventh-century commentary that
names what the loosed mouth first says: \latin{Credo in Deum Patrem
omnipotentem}.

\subsection*{2. Four texts stop where the human contribution would begin}

This is the formulary's most demonstrable coherence, and it is a coherence of
boundaries rather than of themes. Four appointed texts break off at the same
kind of place, and in three of them the breaking is something the liturgy did
to a source that continues.

The Epistle ends at \latin{et grátia eius in me vácua non fuit}. The verse goes
on: \latin{sed abundántius illis ómnibus laborávi: non ego autem, sed grátia
Dei mecum}. Those unread words are the foundation of the entire classical Latin
argument about grace and freedom. Augustine's proof of free will in
\work{De gratia et libero arbitrio} 5.12 is the labour clause, balanced by
\latin{nec gratia Dei sola, nec ipse solus, sed gratia Dei cum illo}. Gregory's
prevenient-and-subsequent scheme in \work{Moralia} 16.25 hangs on one word:
\latin{Non enim diceret \textbf{mecum} si cum praeveniente gratia subsequens
liberum arbitrium non haberet}. Aquinas derives \latin{gratia cooperans} from
the same clause. Chrysostom's reading of the whole passage as humility ``with
reserve'' needs the boast in order to have something retracted. The lection
ends one clause before all of it.

The Collect's second petition asks God to add \latin{quod orátio non
præsúmit} --- and the subject of the verb is \latin{orátio}. It is the prayer,
not the person praying, that does not presume; the 1861 hand missal's ``we dare
not presume to ask'' is a different and more comfortable sentence, and the
divergence is declared where the English is printed. The Communion antiphon
trims Proverbs 3:9 of its imperative \latin{da ei}, so that of the four verbs
in the appointed text three are God's (\latin{implebúntur},
\latin{redundábunt}) or the singer's honouring, and the act of donation has
gone. And the Gospel's protagonist does nothing at all: he is brought, taken
aside, touched with fingers and spittle, sighed over and opened, and the first
thing said of his speech is that it was right.

Two conclusions, one negative and one positive. Negatively, no guide may say
that the Fathers expound this lesson's doctrine of grace: their expositions
quote past its last word, and the guide that omits to say so has quietly
attributed to them a comment on a text the Church does not read here.
Positively, the truncation is old --- the same mid-verse ending stands in the
Venice Missal of 1570 and the Pustet Missal of 1862 --- so it is the received
lectionary boundary of the rite and not a 1962 decision. What the Church
actually reads on this Sunday is a statement about the origin of a man and a
ministry, not about the mechanics of cooperation. That it cannot be turned into
a \textit{denial} of cooperation follows from the same evidence: the authors
whose clause is missing are the ones who teach it.

\subsection*{3. One reading, and the witnesses who refuse it}

Five elements of this Mass have a resurrection reading in the checked corpus,
and the convergence is striking enough that it must be stated carefully. On the
Introit's last clause --- power and strength given to his people --- Hilary
speaks of those whom God will raise as immortal, and Augustine proves the point
from 1 Corinthians 15:43 and Philippians 3:10, arriving at strength ``by which
death the enemy shall be destroyed''; Bellarmine follows him. On the Gradual,
Augustine's gloss is as flat as a gloss can be: \latin{Et refloruit caro mea:
id est, et resurrexit caro mea}. Cassiodorus gives the prefix a history, the
flesh flowering in the virginal conception and re-flowering in the
Resurrection; Lombard transmits both readings correctly tagged; Aquinas extends
them across Adam and Christ; Bellarmine turns them into description. On the
Offertory the psalm's title does the work: Jerome in one line
(\latin{Dedicatio domus David, resurrectio Salvatoris intellegitur}), Augustine
in two expositions, Cassiodorus almost verbatim after him, Aquinas fusing body
and Church into a two-stage dedication. The Epistle is the resurrection
kerygma itself. And Schuster reads the Postcommunion's help to mind and body as
the pledge of the final resurrection.

Now the correction. Not one of those witnesses is commenting on this Sunday;
each is expounding a psalm, a letter or a prayer on its own account. And the
convergence has real dissenters inside the checked corpus. Cassiodorus, who
supplied the Gradual's most influential development, refuses the resurrection
reading at the Introit's last clause and moralises it into patience and faith,
keeping eschatology only as \latin{praemia}. Theodoret refuses it at the
Gradual, hearing a hunted man whose spirits and body have recovered together.
At the Alleluia the split is sharper still and runs the other way: Augustine
makes the psalm's instruments an exchange of spiritual for carnal goods and
Cassiodorus makes the drum a mortified body, while Theodoret keeps the
orchestra literal on the authority of Chronicles and Bellarmine sides with
Theodoret. Four elements, one direction, and three named refusals. The honest
statement is that a reader arriving with the Latin tradition in hand will hear
a resurrection Mass here; the texts supply that reading only through their
reception, and the reception is not unanimous.

\subsection*{4. The body is the instrument in the Gospel and the beneficiary in the orations}

Read across the formulary, this Mass is unusually physical, and the physicality
is doctrinally load-bearing at both ends.

In the Gospel the healing is done with parts of a body. The Greek Mark-catena
tradition that the \work{Catena aurea} labels ``Chrysostomus'' --- in fact the
Victor-of-Antioch material, as the 1842 Oxford editors note --- makes the
assumed body itself the instrument: it was \latin{divina virtute ditatum
\dots\ corpus divinitati unitum}, and in it Christ ``demonstrated the perfection
of human nature.'' Theophylact presses the same point through the least dignified
detail in the story: the spitting was done \latin{ut ostenderet omnes partes
sanctae carnis suae sanctas ac divinas esse}. Bede reads the same gestures
sacramentally in the other direction --- the fingers are the Holy Spirit, on the
warrant of Luke 11:20 read with Matthew 12:28, and the spittle is wisdom
touching the mouth of a catechumen. And the sigh divides them: for Bede it is
\latin{exemplum gemendi}, an example given because Christ needs nothing; for
Victor it is Christ taking our cause on himself; Theophylact will not choose,
hedging twice; a Lapide has it both ways. Augustine, the one Latin master who
touches this pericope at all, ignores the gestures entirely and asks only why
silence was commanded of people who could not keep it.

The orations then make the body the beneficiary. The Secret asks that one gift
be \latin{tibi munus accéptum} and \latin{nostræ fragilitátis subsídium} ---
directed to God, and directed at a weakness admitted in the same breath;
Schuster fixes \latin{servitútem} as the sacerdotal ministry itself, so what
God is asked to look at is the act being performed as the words are said. The
Communion promises filled barns and overflowing presses, which is material
increase promised at the moment of sacramental reception --- and Bede's three
levels keep that from becoming crude, since his middle level makes the
firstfruits the firstfruits of every good work, honoured to God by crediting it
to grace. The Postcommunion refuses to divide the beneficiary at all:
\latin{subsídium mentis et córporis}, and then \latin{in utróque salváti},
resuming exactly those two. Its last verb is \latin{gloriémur}, a verb of
boasting, which the 1861 English renders ``enjoy the full effect'' --- keeping
the petition and losing the boast. Schuster's conclusion, that the body's
contact with the Body of Jesus confers the grace of the final resurrection, is
the point at which this thread rejoins the third: the instrument in the Gospel
and the beneficiary in the orations are the same kind of thing.

\subsection*{5. There is no composer, and the transmission proves it}

Every synthesis above is a statement about the book as it stands, because the
transmission evidence forbids anything stronger. The three orations are the only
part of this Mass that can be shown to have been assembled by anybody: they
stand together as one mass-set in the Old Gelasian (Wilson, p.~228), among the
sixteen unassigned Sunday masses of Book III, tied to no numbered Sunday at
all. Wilson's own footnote records the Frankish Gelasians assigning the set to
the thirteenth Sunday after Pentecost and Pamelius's Gregorian to the twelfth;
the Hadrianum (Wilson, p.~172) files it at \latin{DOMINICA .XII. POST
PENTECOSTEN}, the Gregorian series running one Sunday ahead of 1962 through
this stretch. Two of the three prayers are older than the set: the Secret
stands in the Veronese collection inside an October mass against drought, whose
preface begs for rain, and the Postcommunion in a July mass, with a
\latin{de} that every later book drops. The Collect is absent from the Veronese
book altogether --- a bounded negative established by searching the whole of
Feltoe's text for five distinctive phrases --- so its oldest verified witness is
the Gelasian.

The chants tell a different story again, and the two stories do not meet. The
Ottobonianus margins reported in Wilson's apparatus cue this Sunday's antiphon
set --- \latin{Deus in loco sancto suo}, \latin{Exaltabo te Domine},
\latin{Honora Dominum} --- against \latin{Dominica X post Pentecosten}, two
masses before the orations they now accompany, while Gregorian Dominica XII
carries the chants that 1962 sings on its Thirteenth Sunday. Whatever pairs the
\latin{Honora Dominum} Communion with the \latin{abundantia pietatis} Collect,
it is not a Carolingian decision about this Sunday; it is the later realignment
of two cycles that had been moving independently. Nothing in this guide infers
an intention from the result.

Two footnotes to the transmission, both about texts that outlived their Mass.
The Collect is used verbatim, contracted \latin{adícias} and all, as the Collect
of \latin{Dominica XXVII per annum} in the postconciliar \work{Missale
Romanum}; the reuse was verified in two independent unofficial Latin
transcriptions, no free scan of the official \latin{editio typica tertia} was
located, and the guide reports it at exactly that strength. And the Secret's
variants are printed rather than harmonised: \latin{fragilitati} in the
Veronese and in the Hadrianum's base text, \latin{fragilitatis} in the Gelasian
manuscripts and in Ottobonianus 313, \latin{ad nostram servitutem} in the
Venice Missal of 1570. The 1962 book chose among these. It did not invent them,
and it left no record of choosing.
